%==============================================================================
% Section 2: EC-Weyl Coupling Structure
%==============================================================================
\section{EC--Weyl Coupling Structure}
\label{sec:ecweyl}

This section introduces the basic structure of the Einstein--Cartan connection
and determines the dependence of the EC--Weyl scalar $C^2_\EC$ on the torsion
modes.  The main results are \textbf{Theorem~1 (AX/VT dropout)},
\textbf{Theorem~2 (Palatini protection, EC)}, and \textbf{Theorem~3
($\delta$-sector null theorem)}.  The first two give the selection rule of the
EC--Weyl coupling and the non-propagation of torsion; the third establishes
that no new homogeneous cubic channel appears in the auxiliary MIXING sector.

\subsection{EC Connection and Contortion}
\label{sec:ecweyl:connection}

In Einstein--Cartan theory the connection $\omega^a{}_b$ is treated as a
variable independent of the metric.  Its difference from the Levi--Civita (LC)
connection $\Gamma^a{}_{bc}$ is parametrised by the contortion $K^a{}_{bc}$:
\begin{equation}
  \omega_\EC^a{}_{bc} = \Gamma^a{}_{bc} + K^a{}_{bc}.
\end{equation}
The contortion is determined from the antisymmetric torsion tensor
$T^a{}_{bc} = 2\omega^a{}_{[bc]}$ by
\begin{equation}
  K_{abc} = \frac{1}{2}\bigl(T_{abc} + T_{bca} - T_{cab}\bigr),
  \qquad
  K_{abc} = -K_{acb}.
\end{equation}

For the homogeneous minisuperspace torsion ansatz we adopt three modes:
\begin{itemize}
  \item \textbf{AX mode}: axial scalar $\eta$ (traceless antisymmetric part of
        torsion), with $V=0$ assumed.
  \item \textbf{VT mode}: vector-trace scalar $V$ (trace part of torsion), with
        $\eta=0$ assumed.
  \item \textbf{MX mode}: coexistence of AX and VT, i.e.\ $V\neq 0$ and
        $\eta\neq 0$.
\end{itemize}
The action is
\begin{equation}
  S = \int\left[\frac{R_\EC}{2\kappa^2}
      + \theta\cdot\mathcal{L}_\NY
      + \alpha\cdot C^2_\EC\right]\sqrt{g}\,\dd^4x,
\end{equation}
where $R_\EC$ is the EC Ricci scalar, $\mathcal{L}_\NY$ is the Nieh--Yan
density \cite{NiehYan1982}, and $C^2_\EC$ is the EC Weyl scalar
\cite{Woodard2015}.  The coupling constants are $\alpha\in\mathbb{R}$ (Weyl)
and $\theta\in\mathbb{R}$ (NY).  The effective potential
$\Veff(r,\eta,V;\alpha,\theta)$ is computed explicitly for each topology (script:
\texttt{proofs/ax\_vt\_dropout\_proof.py}).

\subsection{Theorem~1: AX/VT Dropout}
\label{sec:ecweyl:thm1}

\begin{theorem}[AX/VT dropout]\label{thm:dropout}
  In all of $\Sthree\times\Sone$, $\Tthree\times\Sone$, and
  $\Nilthree\times\Sone$, the following hold as exact symbolic identities
  verified by SymPy\textup{:}
  \begin{enumerate}
    \item \textbf{AX mode} ($V=0$, $\eta\neq 0$):
          $C^2_\EC = C^2_\LC$ (no EC correction).
    \item \textbf{VT mode} ($\eta=0$, $V\neq 0$):
          $C^2_\EC = C^2_\LC$ (no EC correction).
  \end{enumerate}
  In AX and VT backgrounds, the EC--Weyl coupling $\alpha C^2_\EC$ coincides
  exactly with the LC case and no $\alpha$-dependent new contribution appears.
\end{theorem}

\begin{proof}[Proof outline]
  Expanding $C^2_\EC - C^2_\LC$ symbolically for each topology and factoring,
  every EC correction term carries $V\eta$ as a common factor.  Therefore in
  the AX mode ($V=0$) the correction vanishes exactly, and in the VT mode
  ($\eta=0$) it likewise vanishes (SymPy exact zero in all cases).
\end{proof}

\paragraph{Topology-specific forms.}
For $\Sthree\times\Sone$ and $\Tthree\times\Sone$:
\begin{equation}
  C^2_\LC(\Sthree, \Tthree) = 0 \quad (\text{isotropic point, AX/VT background}).
\end{equation}
AX dropout then confirms $C^2_\EC = 0$.

For $\Nilthree\times\Sone$:
\begin{equation}
  C^2_\LC(\Nilthree) = \frac{4}{3R^4} \neq 0
  \quad (\text{non-conformal flatness}),
\end{equation}
and AX dropout gives $C^2_\EC = 4/(3R^4) = C^2_\LC$ (no EC correction).

\paragraph{$\Tthree$ squash+shear null test.}
For $\Tthree$, $C^2_\EC = 0$ holds for all $\varepsilon, s$ because
$C^a_{bc}=0$ identically.  These six cases (3 topologies $\times$ 2 backgrounds)
are all confirmed as exact zero by SymPy (scripts: \texttt{proofs/ax\_vt\_dropout\_proof.py},
\texttt{proofs/t3\_flatness\_null\_test.py}).

\subsection{MX Mode: EC-Specific Weyl Coupling}
\label{sec:ecweyl:mx}

As the converse of Theorem~\ref{thm:dropout}, in the MX background
($V\neq 0,\eta\neq 0$) an EC-specific Weyl coupling appears.

For $\Sthree\times\Sone$ and $\Tthree\times\Sone$:
\begin{equation}
  C^2_\EC = C^2_\LC + \frac{16V^2\eta^2}{3r^2}.
\end{equation}
For $\Nilthree\times\Sone$:
\begin{equation}
  C^2_\EC = C^2_\LC + \frac{16V^2\eta^2}{3R^2}
           = \frac{4}{3R^4} + \frac{16V^2\eta^2}{3R^2}.
\end{equation}
In both cases the additional term is proportional to $(V\eta)^2$, so the
torsion mixing product modifies the effective potential via the EC--Weyl
coupling:
\begin{equation}
  \Delta\Veff^\EC = \alpha\cdot\frac{16V^2\eta^2}{3r^2}\cdot\Vol(M^3\times\Sone).
\end{equation}
This EC-specific term is the source of effective-potential deformation in the
MX background.  Note, however, that the $\Nilthree$ EC slice minimum discussed
in \S\ref{sec:nil3} originates from $C^2_\LC(\Nilthree)\neq 0$ on the
$\eta=V=0$ section, which is independent of this MX term.

\paragraph{Physical interpretation.}
AX/VT dropout means that torsion alone does not generate new contributions to
the Weyl tensor.  The essence of the EC--Weyl interaction lies in the
``mixing product'' $V\times\eta$ of torsion, understood geometrically as the
mechanism by which an asymmetric combination of contortion activates Weyl
components.  This selection rule is universal across all three topologies and
follows from the general algebraic structure of the EC connection.

\subsection{Theorem~2: Palatini Protection (EC Version)}
\label{sec:ecweyl:thm2}

\begin{theorem}[Palatini protection, EC]\label{thm:palatini}
  Under the Einstein--Cartan connection,
  \begin{equation}
    G^{\EC}_{\eta\eta} = G^{\EC}_{VV} = 0
  \end{equation}
  holds exactly, where $G^{\EC}_{AB} =
  \partial^2\mathcal{L}_\EC/\partial(\partial_z\Phi^A)\partial(\partial_z\Phi^B)$
  is the field-space metric (kinetic metric) of the minisuperspace.
\end{theorem}

\begin{proof}[Proof outline (velocity method)]
  The EC Weyl scalar in the MX background,
  $C^2_\EC(\Sthree, \Tthree, \text{MX}) = 16V^2\eta^2/(3r^2)$,
  is an algebraic function of $(r,\eta,V)$ and contains no velocity fields
  $\partial_z\eta$, $\partial_z V$.  Therefore, computing $G_{\eta\eta}$,
  $G_{VV}$ by the velocity method yields only terms that vanish in the
  $\Delta z\to 0$ limit, and no kinetic term for torsion is generated.
  Full symbolic verification for each topology is given in script
  \texttt{proofs/palatini\_protection\_proof.py}.
\end{proof}

\paragraph{Invariance of $r$-stiffness.}

\begin{table}[H]
\centering
\begin{tabular}{lccc}
\toprule
Topology & $G_{rr}$ (LC) & $G_{rr}$ (EC) & Difference \\
\midrule
$\Sthree$ & 59.22 & 59.22 & 0 \\
$\Tthree$ & 4675.6 & 4675.6 & $\approx 0$ \\
$\Nilthree$ & 5195.2 & 5195.2 & $\approx 0$ \\
\bottomrule
\end{tabular}
\caption{Kinetic metric $G_{rr}$ for the radial mode: LC vs.\ EC.  The
  graviton kinetic structure is unchanged under the EC extension.}
\label{tab:Grr}
\end{table}

The kinetic structure of the spin-2 graviton is also invariant under EC.

\subsection{Theorem~3: $\delta$-Sector Null Theorem}
\label{sec:ecweyl:thm3}

In the EC extension, the new nonlinear coupling
\begin{equation}
  \frac{\partial^3 V_\EC}{\partial\alpha\,\partial\eta\,\partial V} \neq 0
\end{equation}
appears in the physical $(\eta,V)$ channel (\S\ref{sec:s3:ec}).  A separate
question is whether an independent cubic channel also arises in the auxiliary
MIXING variable $\delta_k$.  We define
\begin{equation}
  C_\delta(M^3;\text{mode}) :=
  \left.\frac{\partial^3\Veff}{\partial\delta_0\,\partial\delta_1\,\partial\delta_2}
  \right|_{\delta_0=\delta_1=\delta_2=0},
\end{equation}
where $M^3 = \Sthree, \Tthree, \Nilthree$ and mode $=$ AX/VT/MX.

\begin{theorem}[$\delta$-sector null theorem]\label{thm:delta}
  For all three topologies and all torsion backgrounds \textup{(}AX, VT,
  MX\textup{)},
  \begin{equation}
    C_\delta(M^3;\text{mode}) = 0
    \qquad\text{and}\qquad
    \frac{\partial C_\delta}{\partial\alpha} = 0
  \end{equation}
  hold as exact SymPy results.  The EC extension generates no new homogeneous
  cubic channel in the auxiliary $\delta$-sector.
\end{theorem}

\begin{proof}[Proof outline]
  With \texttt{fiber\_mode = MIXING} activated in the unified engine, we
  construct $\Veff$ for each topology and torsion background, then compute the
  third derivative $\partial_{\delta_0}\partial_{\delta_1}\partial_{\delta_2}\Veff|_{\delta=0}$
  by direct symbolic differentiation.  All nine cases give exact zero
  coefficients, from which $\partial_\alpha C_\delta = 0$ follows automatically.
  The key point is not that EC has no new effects, but that the new cubic
  structure activated by EC is localised in the physical $\eta$-$V$ channel
  rather than the auxiliary $\delta$-sector.
\end{proof}

\paragraph{Topology-wise consequences.}
\begin{itemize}
  \item $\Sthree$: The formal MIXING sector is nontrivially available, yet
        $C_\delta(\Sthree)=0$.  Hence the $\alpha$-cubic of \S\ref{sec:s3:ec}
        appears only in the $\eta$-$V$ sector.
  \item $\Tthree$: Flatness trivialises most of the mode dictionary.  Even the
        formal MIXING test gives $C_\delta(\Tthree)=0$.
  \item $\Nilthree$: Despite the non-conformal flatness and the EC false
        vacuum, $C_\delta(\Nilthree)=0$.  The EC-specific physics of $\Nilthree$
        resides in the background Weyl and the $(r,\eta,V)$ sector.
\end{itemize}
Detailed derivations and the exact check of all nine cases are given in
Appendix~\ref{app:D} and script \texttt{proofs/c\_delta\_null\_proof.py}.
